\section{引言}

\subsection{辛奇点之屋}
\label{sec:house}

有人把数学研究比作在一间漆黑的屋子里摸索电灯的开关\footnote{「也许最能描述我做数学的体验的，是走进一座漆黑的宅邸。你走进第一个房间，里面漆黑一片，一片彻底的黑暗。你跌跌撞撞，不断撞上家具；渐渐地，你知道了每件家具都在什么位置；最后，大约过了半年，你摸到了电灯开关。你打开开关，刹那间一切都被照亮，你能清清楚楚地看到自己刚才究竟身在何处。」——Andrew Wiles}。在另一些情形下，它又像是在白天穿过同一座宅子——所有家具都看得见，但你仍可以翻箱倒柜，搜寻小块的珍宝。这样的寻宝固然惬意（对小leg也友好得多），然而发现我们从未见过的房间，可能带来更大的收获。在某些领域，这不过是沿着走廊走一趟的事，困难的只是知道该推开哪一扇门。而更令人兴奋的，是在两个我们自以为早已熟悉的房间之间发现一条密道。

当然，除非你正在玩一局《妙探寻凶》（Clue），密道并不容易找到。你总不能把墙拆开，还指望密道就在那里。然而，在20世纪末、21世纪初，数学家找到了这类密道的一个惊人来源：量子场论（quantum field theory，QFT）。

QFT中所谓的「对偶」，常常把事先看来毫无关联的数学对象联系起来。例如，\textbf{（二维）镜像对称}表明，代数几何学家与辛几何学家其实一直住在同一座宅子里，只不过他们之间的那条通道至今仍然相当昏暗，走起来也比我们希望的要费力。

遗憾的是，要在数学里用上这些对偶，并不是请一位物理学家戴上他的X光透视眼镜（x-ray specs）过来就能了事；它更像是收到一套残缺不全、饱经风霜的蓝图，上面的文字或许还是一种陌生的语言，只能暗示你该往哪里看。尽管如此，我们还是得到了一些非常有趣的提示。

对表示论学家来说，所有宅邸中最辉煌、也勘探得最充分的一座，是单李代数之屋；它已有一百多年的历史，却仍有许多幽深的角落藏着迷人的惊喜。它还有一处看上去相当无害的小过道，连通不同的房间，叫作Langlands对偶。毕竟，那不过是把Cartan矩阵转置一下；而我们大多数人若不查书，本来就分不清Cartan矩阵和它的转置。Langlands纲领揭示了这一简单操作深不可测的内涵。

这座庄园还陆续添了许多新的侧翼：李超代数、特征 $p$ 的代数群表示、箭图表示、量子群、范畴化，等等。它们虽然五花八门，却都依托于Dynkin图这同一个底层框架。但近年来，研究者找到了又一种推广，它的意义更类似于发现许多新的Dynkin图系列：\textbf{辛消解}与\textbf{辛奇点}的世界。用一句常被引用、通常归在Okounkov名下的妙语（bon mot）来说：「辛奇点是21世纪的李代数。」

\begin{yizhu}[辛消解与辛奇点]
\textbf{辛消解}指仿射簇 $X_0$ 的一个消解 $X\to X_0$，其中 $X$ 是带辛形式的光滑簇；\textbf{辛奇点}则指这样的仿射簇：其光滑轨迹上配有辛形式，且该形式在奇点处性态良好——$X_0$ 正是它的特例。二者被视为李代数在21世纪的类比物：辛消解扮演Springer消解的角色，而其「泛包络代数」扮演 $U(\mathfrak{g})$ 的角色。
\end{yizhu}

关于这处特别的附楼，有趣的结果大约在21世纪之交开始出现，其基础是Kaledin、Bezrukavnikov等人的工作。2007年前后的某一段时间，我的\footnote{本节中的所有代词都取自B.W.的视角。}合作者Tom Braden、Nick Proudfoot、Tony Licata和我注意到了另一条密道的迹象，它在那里连接着一对对房间（即辛消解）。要让这些不同的房间精确对齐、为一条密道腾出空间，需要许多巧合同时发生。然而，我们没能踏进这条通道本身。尽管如此，我们找到了一个非常引人入胜的例子：我们寻找的密道，将把Langlands对偶推广到许多新的情形。

当然，从上面的讨论你就能猜到后来发生了什么。2008年，我在高等研究院（Institute for Advanced Study）作了一场报告，之后Sergei Gukov向我指出：物理学家早已知道，基于一个已知的对偶——\textbf{三维镜像对称}——这些密道应当存在。如前所述，这绝对没有解决我们的全部问题；直到今天，我们当时作出的若干观察仍未得到解释。更一般地说，这个对偶在当时并未被物理学家很好地理解（而且许多问题至今悬而未决），但它至少解释了这样一条通道为何应当存在，并为进一步寻找它提供了依据。

自那次交谈以来的15年间，数学与三维QFT之间的联系取得了巨大的进展。本文的目的是向数学家简短地解释这些进展，以及其背后的一部分QFT。它必然极不完整，这一点令人遗憾，但我们希望它能成为各个年龄段的数学家进一步学习的有用指南。

\subsection{全文安排}

现在我们用稍微精确一点的语言来说明全文的安排。\textbf{辛消解}是由以下两者组成的一对：
\begin{enumerate}
    \item 一个带奇点的仿射簇 $X_0$；以及
    \item 一个带代数辛形式的光滑簇 $X$，它消解 $X_0$ 的奇点。
\end{enumerate}
这个带奇点的仿射簇 $X_0$ 是\textbf{辛奇点}的一个特例；所谓辛奇点，是指这样的带奇点的仿射簇：其光滑轨迹上配有辛形式，且该形式在奇点处性态良好。

辛消解最著名的例子是\textbf{Springer消解}：其中 $X_0$ 是半单李代数 $\mathfrak{g}$ 中幂零元构成的簇，$X$ 是 $\mathfrak{g}$ 的旗簇的余切丛。由这个消解的几何可以重构出 $\mathfrak{g}$。因此，看待单李代数之屋的一种方式是：Springer消解才是单李代数每个房间里的基本对象，而李理论的其他一切方面，都由从各种不同角度审视Springer消解所决定。

于是，单李代数位于一条走廊的一端，走廊上还有许多别的门，通向其他的辛消解与辛奇点。这就引出一个自然的问题：如果我们把别的辛消解当作某个我们从未见过的、新的李代数的Springer消解，那么李代数的任一给定概念是否都能推广过去？例如，每个辛奇点都有一个「泛包络代数」，它是李代数的泛包络代数之推广。

大多数数学家都能上手的两个例子是：
\begin{itemize}
    \item 复射影空间的余切丛 $\XA=T^*\mathbb{CP}^{n-1}$。它可以写成
    \begin{equation*}
    T^*\mathbb{CP}^{n-1} =\{(\ell, \phi)\in \mathbb{CP}^{n-1}\times M_{n\times n}(\C) \mid \phi(\C^n)\subset \ell, \phi(\ell)=\{0\}\}
    \end{equation*}
    到第二个分量的投影是 $\rkone$ 的一个消解，这里 $\rkone$ 表示秩 $\leq 1$ 的 $n\times n$ 矩阵构成的空间。
    这个余切丛带有典范辛形式，从而使该消解成为辛消解。
    \item 循环群 $\Z/n\Z$ 通过矩阵
    \[k\mapsto \begin{bmatrix} \exp(2\pi ik/n) & 0\\0& \exp(-2\pi ik/n)\end{bmatrix}
    \]
    作用在 $\C^2$ 上，并保持其典范辛形式。商 $\C^2/(\Z/n\Z)$ 有唯一的辛消解 $\XB$，其例外纤维是 $n-1$ 个 $\mathbb{CP}^1$ 的并，这些 $\mathbb{CP}^1$ 连成一条链。
\end{itemize}
有同构 $X_2^{(\mathrm{A})}\cong X_2^{(\mathrm{B})}$，但当 $n>2$ 时，这两个簇的维数并不相同。
若考察从这些簇中得出的某些组合信息，就会发现一些耐人寻味的共同点。其中居核心地位的是两个几何对象：
\begin{itemize}
    \item 极大环面 $T^{(*)}$（$*\in \{\mathrm{A},\mathrm{B}\}$）在 $\XS$ 上的作用，它保持辛结构。一个显然的不变量是该环面的不动点集。\footnote{$T^{(\mathrm{A})}$ 是 $PGL_n(\C)$ 中的对角矩阵；$T^{(\mathrm{B})}$ 是 $SL_2(\C)$ 中的对角矩阵模去 $n$ 挠。}
    \item 仿射簇 $X_0$ 有唯一的极小分解，分解为有限多个带诱导辛结构的光滑块，这是幂零矩阵按Jordan型分解的推广。
\end{itemize}
这对簇之间有一些耐人寻味的巧合：
\begin{enumerate}
    \item 有同构
    \[\mathfrak{t}^{(\mathrm{A})}\cong H^2(\XB)\qquad \mathfrak{t}^{(\mathrm{B})}\cong H^2(\XA).\]
    我们可以把这个结论加强：注意这两个空间上几何地定义出的超平面配置是相配的。\footnote{在 $\mathfrak{t}$ 中，取的是使相应向量场的零点集维数发生跳跃的那些向量；在 $H^2(\XB)$ 中，取的是切出同一仿射簇各个不同crepant消解的丰沛锥的那些Mori壁。}
    \item 两个环面作用的不动点个数相同，都是 $n$；这也说明 $\XS$ 的Betti数之和为 $n$。
    \item $X_0^{(\mathrm{A})}$ 与 $X_0^{(\mathrm{B})}$ 上的分层片数相同，都是 $2$。\footnote{光滑轨迹是其中一片；在两种情形下，另一片都是单点。}
\end{enumerate}
人们很容易把这些当作无关紧要而一笑置之，但它们是一个更丰富现象的数值体现。也就是说，
\begin{enumerate}
    \item[4.] $\XS$ 的「泛包络代数」有一类特殊的表示，我们称之为「范畴 $\mathcal{O}$」（见 \cite[\S 3]{BLPWgco}），而 $\XA$ 与 $\XB$ 的范畴 $\mathcal{O}$ 互为\textbf{Koszul对偶}；一个范畴中投射模之间的同态，刻画了另一个范畴中单模之间的扩张。
\end{enumerate}
我们不应把这些「巧合」一笑置之的另一个理由是：同样的命题 1.--4. 对许多对辛奇点都成立，\cite[\S 9]{BLPWgco} 中对此有讨论。这些例子包括所有有限型与仿射型A的箭图簇，以及光滑的超环面簇。有些例子是自对偶的：
\begin{itemize}
    \item $Y_n^{(\mathrm{A})}=Y_n^{(\mathrm{B})}=T^*\operatorname{Fl}_n$，即 $\C^n$ 中完全旗构成的簇的余切丛。
    \item $Z_n^{(\mathrm{A})}=Z_n^{(\mathrm{B})}=\operatorname{Hilb}^{n}(\C^2)$，即 $\C^2$ 中 $n$ 个点的Hilbert概形。
\end{itemize}
对第 3 条作适当修改\footnote{此时各层是幂零元的伴随轨道，其数目对 $B_n$ 型与 $C_n$ 型并不相同。若只考虑\textbf{特殊轨道}，则数目相同，可以恢复一一对应。}之后，它还包括Langlands对偶的李代数对的Springer消解。

\begin{yizhu}[超环面簇]
\textbf{超环面簇}（hypertoric variety，又称环面超凯勒簇）由环面在超凯勒向量空间上的超凯勒商得到，可以看作箭图簇的环面类比；它自带辛结构，因而也是辛奇点的来源之一。
\end{yizhu}

这种作用于辛奇点及其消解之上的神秘对偶，由于把两个看似无关的辛簇联系了起来，得到了一个名字：「辛对偶」。

\begin{ques}\label{ques:duality}
是否存在一条底层原理，能够解释前面第 1–4 条陈述，也就是说，能够解释这些簇对之间的辛对偶？
\end{ques}

如前所述，三维QFT方面的工作提示，这个问题的答案离「是」比离「否」更近。本文的目标，就是解释这一点的基本道理，以及它告诉了我们关于数学的什么。

我们可以把它拆成两个子问题：
\begin{enumerate}
    \item[Q1.] 三维 $\mathcal{N}=4$ 超对称量子场论及其拓扑扭转是什么？
    \item[Q2.] 它们与辛对偶有什么关系？
\end{enumerate}
在第\ref{sec:physical-origin}节中，我们将回答这两个问题，这里先作简要概括。

首先，每个三维拓扑量子场论（topological quantum field theory，TQFT）都给出一个泊松代数（关于这一构造的更多讨论见 \cite{BBBDN}）。在许多情形下，这个环是某个辛奇点 $X_0$ 的坐标环，而上面讨论的所有例子都可以这样构造出来。

给定一个QFT，选定一种拓扑扭转，便得到一个TQFT。事实上，对一个三维 $\mathcal{N}=4$ 理论 $\mathcal{T}$，这样的选择有两种，分别称为 $A$ 扭转与 $B$ 扭转。因此，每个三维理论都给出两个辛奇点 $\fM_A(\mathcal{T})$ 与 $\fM_B(\mathcal{T})$，分别称为该理论的\textbf{库仑分支}与\textbf{希格斯分支}。

那些辛簇对 $X^{(\mathrm{A})}$ 与 $X^{(\mathrm{B})}$（类似地还有 $Y$、$Z$ 等），结果都是同一个理论 $\mathcal{T}$ 的库仑分支与希格斯分支。于是命题 1.--4. 就可以用第\ref{sec:house}节提到的那个物理对偶来理解，它叫作「三维镜像对称」。

这是一个非常庞大的题目；由于篇幅和参考文献数量的限制，我们将集中于它与奇点辛消解的关系 \cite{beauvilleSymplecticSingularities2000}。特别地，对于物理学中关于此题的悠久而丰富的文献，我们只会相对简略地带过；读者可以从 \cite{bullimoreBoundariesMirror2016} 的引言出发找到相关文献，从奠定这一理论基石的 Intrilligator--Seiberg \cite{intriligatorMirrorSymmetry1996} 与 Hanany--Witten \cite{hananyTypeIIB1997} 的原始工作开始。

正如数学家所熟悉的二维镜像对称提示我们：复流形与（带额外结构的）辛流形成对出现，而它们之间的关系起初很难察觉；三维镜像对称则改写了我们对问题\ref{ques:duality}的回答：一个理论的库仑分支也可以看作它的对偶理论的希格斯分支，即 $\fM_A(\mathcal{T})=\fM_B(\mathcal{T}^{\vee})$。因此，我们也可以把那些对偶的辛簇对描述为对偶理论 $(\fM_B(\mathcal{T}), \fM_B(\mathcal{T}^{\vee}) )$ 的希格斯分支。

这个答案并不像我们希望的那样完整，因为我们无法把三维QFT构造为严格的数学对象。我们只能对某些类型理论的某些方面做严格的数学处理，其中最重要的是\textbf{线性规范理论}。在这些情形下，希格斯分支与库仑分支都有数学定义，因而我们可以证明关于它们的数学结论。

在第\ref{sec:Higgs-and-Coulomb}节中，我们将回顾线性规范理论情形下希格斯分支与库仑分支的上述构造。前一种构造数学家已经知道几十年了 \cite{hitchinHyperkahlerMetrics1987}；而库仑分支的构造在2015年出现时，连物理学家都感到意外 \cite{BFN}，它也是我们此后取得进展的关键。

这些簇是数学与物理交叉的一个迅速发展中的研究领域的拱心石。特别地，它们为理解三维理论的镜像对称指明了道路：这种镜像对称不仅是数学家已知的镜像对称（即二维镜像对称）的对应物，还丰富了几何Langlands纲领（后者来自四维理论的对偶）。

我们将在第\ref{sec:advanced}节结束本文，简要讨论当前与未来一些有趣的研究方向，为有兴趣的读者指引下一步该往哪里走。

\section*{致谢}

帮助过我们学习这些材料的同事与合作者，包括对本文提出反馈的诸位，多到几乎无法一一列举。Justin Hilburn尤其值得感谢，他与我们有过许多讨论，并就这些题目贡献了许多洞见；同样值得感谢的还有Tudor Dimofte、Davide Gaiotto、Kevin Costello、David Ben-Zvi、Nick Proudfoot、Tom Braden、Anthony Licata、Contantin Teleman、Andrei Okounkov、Mina Aganagi\'c、Richard Rimany\'i、Hiraku Nakajima与Sam Raskin。

这项研究还得到Perimeter Institute for Theoretical Physics的支持：B.W.在那里度过了六年，P.Y.也曾多次访问，此外还有两次「QFT for Mathematicians」会议，它们对本文的形成至关重要。Perimeter Institute的研究由加拿大政府（通过创新、科学与经济发展部）以及安大略省政府（通过研究创新部）支持。本工作还得到加拿大自然科学与工程研究理事会（Natural Sciences and Engineering Research Council of Canada）发现基金RGPIN-2018-03974、首尔国立大学新聘教员启动基金，以及韩国政府（MSIT）资助的韩国国家研究基金会（NRF）项目No. 2022R1F1A107114212的支持。

% [原文 p1–p3]
